\documentclass{article}

% Language setting
% Replace `english' with e.g. `spanish' to change the document language
\usepackage[english]{babel}

% Set page size and margins
% Replace `letterpaper' with `a4paper' for UK/EU standard size
\usepackage[
      letterpaper,
      top=2cm,
      bottom=2cm,
      left=3cm,
      right=3cm,
      marginparwidth=1.75cm
]{geometry}

% Useful packages
\usepackage{amsmath}
\usepackage{graphicx}
\usepackage[colorlinks=true, allcolors=blue]{hyperref}

\title{Draft - Quantum-Safe Messaging Protocol}
\author{Noella Alexandra Mabelle}

\begin{document}
      \maketitle
      \section{Whatsapp Signal Protocol}
      WhatsApp utilise le Signal Protocol — X3DH pour établir une connexion sécurisée,
      combiné au Double Ratchet pour changer les clés après chaque message. Même
      si une clé est compromise, tous les messages passés et futurs restent
      protégés.
      \begin{itemize}
            \item Every user has three type of keys on the server of Whatsapp : une
                  identity key (long terme, pour verifier l’identité), une signed
                  prekey (semi-long terme, pour initier une session même si long
                  terme), et des one-time prekeys (usage unique).

            \item Use of X3DH to send a message $SS1 = DH(IK_{A}, SPK_{B}) ← iden
                  tité A lice + signed prekey Bob SS2 = DH(EK_{A}, IK_{B}) ← éphé
                  mère Alice + identi té Bob SS3 = DH(EK_{A}, SPK_{B}) ← éphémère
                  Alice + signed prekey Bob SS4 = DH(EK_{A}, OPK_{B}) ← éphémère
                  Alice + one-time prekey Bob Secret final = H KDF(SS1 || SS2 ||
                  SS3 || SS4)$

            \item Double ratchet, X3DH (the signal protocol)

            \item HANDSHAKE : X3DH, 4X ECDH (X25519)

            \item MESSAGES (session) : Double Ratchet, Symmetric + DH Ratchet (ECDH)

            \item AUTHENTICATION : ECDSA (Ed25519)

            \item SYMETRIC : AES-256-CBC + HMAC-SHA256
      \end{itemize}

      \section{Quantum safe algorithm}
      Implement a point-to-point messaging prototype inspired by WhatsApp
      covering handshake, authentication, session management, and forward secrecy.
      \begin{itemize}
            \item Following NIST 2024 standards to be fully quantum safe

            \item Simplified version (no groups, no multi-devices)

            \item Quantum safe key exchange: ML-KEM

            \item Quantum safe authentication: ML-DSA - AES-256 is kept for
                  symmetric encryption

            \item Simplified ratchet

            \item To evaluate the practical implications of post-quantum messaging
                  : o MIM simulation : what an attacker can capture on classical
                  vs quantum-safe protocols o Comparative Benchmark of key sizes,
                  handshake latency and signature sizes for classical and post-quantum.
      \end{itemize}

      \subsection{Definition of the algorithms used}
      This protocol combines post-quantum public-key primitives with classical
      symmetric cryptography. The objective is to replace the parts of a
      messaging protocol that are vulnerable to quantum computers, while keeping
      efficient symmetric protection for actual message contents.

      \subsubsection{ML-KEM}
      ML-KEM is the \textbf{Module-Lattice-Based Key Encapsulation Mechanism}
      standardized by NIST in FIPS 203. It is the post-quantum replacement for
      classical key exchange mechanisms such as Diffie--Hellman or X25519.

      Its role in the protocol is to let two parties derive a shared secret over
      an untrusted network without exposing that secret to an eavesdropper. At a
      high level, ML-KEM works in three steps:
      \begin{enumerate}
            \item \textbf{Key generation:} a receiver generates a public key and
                  a secret key.
            \item \textbf{Encapsulation:} the sender uses the receiver's public
                  key to generate a ciphertext and a shared secret.
            \item \textbf{Decapsulation:} the receiver uses the secret key to
                  recover the same shared secret from the ciphertext.
      \end{enumerate}

      In contrast with classical Diffie--Hellman, the security of ML-KEM is
      based on hard lattice problems rather than elliptic-curve discrete
      logarithms. This makes it resistant to known quantum attacks such as
      Shor's algorithm. In this project, ML-KEM is used during the handshake and
      during periodic rekeying in the simplified ratchet.

      \subsubsection{ML-DSA}
      ML-DSA is the \textbf{Module-Lattice-Based Digital Signature Algorithm}
      standardized by NIST in FIPS 204. It is the post-quantum replacement for
      classical signature schemes such as Ed25519 or RSA signatures.

      Its role in the protocol is authentication. It proves that a public key,
      handshake message, or ratchet update was really produced by the claimed
      user and not by an attacker or by the server. At a high level, ML-DSA
      works as follows:
      \begin{enumerate}
            \item \textbf{Key generation:} a user creates a signing key and a
                  verification key.
            \item \textbf{Signing:} the sender computes a signature on a message
                  or transcript using the signing key.
            \item \textbf{Verification:} the receiver checks that signature using
                  the public verification key.
      \end{enumerate}

      The security of ML-DSA also comes from structured lattice problems rather
      than factorization or discrete logarithms. In this messaging design, ML-DSA
      signs identity information, signed prekeys, handshake transcripts, and
      optional ratchet public keys. This is what prevents man-in-the-middle key
      substitution attacks.

      \subsubsection{AES-256-GCM}
      AES-256-GCM is a symmetric authenticated encryption algorithm. AES
      (Advanced Encryption Standard) provides confidentiality, and GCM
      (Galois/Counter Mode) adds integrity and authenticity.

      Its role in the protocol is to protect the actual content of messages once
      a shared session key has been established. At a high level:
      \begin{enumerate}
            \item a symmetric message key is derived from the session state,
            \item AES encrypts the plaintext into ciphertext,
            \item GCM computes an authentication tag over both the ciphertext and
                  the additional authenticated data.
      \end{enumerate}

      AES-256 is kept in the post-quantum setting because symmetric encryption is
      much less affected by quantum computing than public-key cryptography. A
      quantum search attack such as Grover's algorithm only gives a quadratic
      speedup, so a 256-bit symmetric key still provides a strong security
      margin.

      \subsubsection{HKDF}
      HKDF is the \textbf{HMAC-based Key Derivation Function}. It is not a
      public-key algorithm: instead, it is used to transform shared secrets into
      structured cryptographic keys.

      Its role in the protocol is to derive:
      \begin{itemize}
            \item the master secret from the handshake secrets,
            \item the root key,
            \item the sending and receiving chain keys,
            \item the per-message encryption keys.
      \end{itemize}

      At a high level, HKDF takes input keying material and expands it into one
      or several pseudorandom outputs that are independent from each other. This
      is important because the protocol should never reuse one raw secret
      directly for several cryptographic purposes.

      \subsubsection{Hash function}
      A cryptographic hash function is used to compute a fixed-size digest of the
      handshake transcript. In practice, the protocol can use SHA-384 or SHA-512
      consistently with the HKDF choice.

      Its role in the protocol is to bind all handshake messages to the final key
      derivation. If an attacker modifies even one field in the transcript, the
      transcript hash changes and the resulting session key will not match. This
      helps protect the integrity of the handshake context.

      \subsubsection{Signed prekeys and ephemeral keys}
      Even though signed prekeys and ephemeral keys are not separate algorithms,
      they are important protocol mechanisms inspired by WhatsApp and Signal.

      \begin{itemize}
            \item A \textbf{signed prekey} is a semi-static ML-KEM public key
                  signed by the user's ML-DSA identity key. It allows a new
                  session to be started without reusing the identity key for key
                  exchange.
            \item An \textbf{ephemeral key} is a fresh temporary ML-KEM key pair
                  created for a single handshake or rekey event. Once used, the
                  private part must be erased.
      \end{itemize}

      These mechanisms improve forward secrecy. Even if a long-term identity key
      is later exposed, previously deleted ephemeral secrets should not be
      recoverable.

      \subsection{Protocol Flow }
      The proposed protocol keeps the WhatsApp design philosophy: the server
      relays public material and encrypted messages, but it must not learn the
      session keys. In the classical Signal design, X3DH and the Double Ratchet
      provide authentication and forward secrecy. In the post-quantum version,
      we replace elliptic-curve primitives with NIST-standardized post-quantum
      mechanisms:
      \begin{itemize}
            \item \textbf{Key establishment:} ML-KEM (FIPS 203).
            \item \textbf{Authentication:} ML-DSA (FIPS 204).
            \item \textbf{Symmetric protection:} AES-256-GCM for confidentiality
                  and integrity.
            \item \textbf{Key derivation:} HKDF-SHA-384 or HKDF-SHA-512.
      \end{itemize}

      This prototype is intentionally simplified: there are no groups, no
      multi-device synchronization, and no advanced skipped-message handling.
      The design only targets point-to-point communication between one device
      owned by Alice and one device owned by Bob.

      \subsubsection{Long-term and semi-static keys}
      Before any communication, each user locally generates the following keys:
      \begin{enumerate}
            \item \textbf{Identity signature key pair}
                  $(IK^{sig}_{priv}, IK^{sig}_{pub})$ using ML-DSA. This is the
                  long-term identity of the user and is used to sign protocol
                  material.
            \item \textbf{Signed prekey KEM pair}
                  $(SPK^{kem}_{priv}, SPK^{kem}_{pub})$ using ML-KEM. This is a
                  medium-term key used to initiate new sessions.
            \item \textbf{Optional one-time prekeys}
                  $(OPK^{kem}_{priv}, OPK^{kem}_{pub})$ if asynchronous messaging
                  is required. For a minimal online prototype, these can be
                  omitted.
      \end{enumerate}

      Bob publishes a \emph{prekey bundle} to the server containing:
      \[
            \{ID_B,\ IK^{sig}_{pub,B},\ SPK^{kem}_{pub,B},\ expiry,\ version,\ \sigma_B\}
      \]
      where
      \[
            \sigma_B = \text{ML-DSA.Sign}(IK^{sig}_{priv,B}, ID_B \parallel SPK^{kem}_{pub,B} \parallel expiry \parallel version).
      \]
      The server stores the bundle, but it cannot forge it because it does not
      know Bob's ML-DSA private identity key.

      \subsubsection{Step 1: Bundle retrieval and verification}
      When Alice wants to start a conversation, she queries the server for Bob's
      prekey bundle. She receives
      \[
            \{IK^{sig}_{pub,B},\ SPK^{kem}_{pub,B},\ \sigma_B\}.
      \]
      Alice verifies $\sigma_B$ with Bob's public identity key. If the signature
      check fails, the protocol aborts immediately. This verification step is
      essential because, without it, the server or a man-in-the-middle attacker
      could replace Bob's public key material.

      \subsubsection{Step 2: Alice sends the handshake initiation}
      Alice generates a fresh ephemeral ML-KEM key pair
      $(EK^{kem}_{priv,A}, EK^{kem}_{pub,A})$ and a fresh nonce $N_A$. She then
      signs the handshake context using her ML-DSA identity key:
      \[
            \sigma_A = \text{ML-DSA.Sign}(IK^{sig}_{priv,A},
            ID_A \parallel ID_B \parallel EK^{kem}_{pub,A} \parallel N_A \parallel context).
      \]
      Alice sends the following \texttt{HANDSHAKE\_INIT} message to Bob:
      \[
            \{version,\ ID_A,\ ID_B,\ EK^{kem}_{pub,A},\ N_A,\ \sigma_A\}.
      \]
      This message does not yet reveal the final session key, but it binds
      Alice's identity to her fresh ephemeral KEM contribution.

      \subsubsection{Step 3: Bob authenticates Alice and replies}
      After receiving \texttt{HANDSHAKE\_INIT}, Bob fetches Alice's public
      identity key and verifies $\sigma_A$. If the verification succeeds, Bob
      generates his own fresh ephemeral ML-KEM key pair
      $(EK^{kem}_{priv,B}, EK^{kem}_{pub,B})$ and a nonce $N_B$.

      Bob then encapsulates to Alice's ephemeral public key:
      \[
            (ct_B,\ ss_B) = \text{ML-KEM.Encaps}(EK^{kem}_{pub,A}).
      \]
      This produces a ciphertext $ct_B$ and a shared secret $ss_B$ that only
      Alice can recover by decapsulation with $EK^{kem}_{priv,A}$.

      To make the exchange mutually contributive, Alice will also encapsulate to
      Bob's ephemeral public key after receiving the reply. Bob signs the full
      transcript:
      \[
            \sigma'_B = \text{ML-DSA.Sign}(IK^{sig}_{priv,B},
            ID_A \parallel ID_B \parallel EK^{kem}_{pub,A} \parallel EK^{kem}_{pub,B}
            \parallel N_A \parallel N_B \parallel ct_B \parallel context).
      \]
      Bob sends the following \texttt{HANDSHAKE\_REPLY} message:
      \[
            \{version,\ ID_B,\ ID_A,\ EK^{kem}_{pub,B},\ ct_B,\ N_B,\ \sigma'_B\}.
      \]

      \subsubsection{Step 4: Alice completes the key establishment}
      Alice verifies Bob's signature $\sigma'_B$, then decapsulates:
      \[
            ss_B = \text{ML-KEM.Decaps}(ct_B, EK^{kem}_{priv,A}).
      \]
      Next, Alice encapsulates to Bob's fresh ephemeral public key:
      \[
            (ct_A,\ ss_A) = \text{ML-KEM.Encaps}(EK^{kem}_{pub,B}).
      \]
      She sends the resulting ciphertext to Bob in a short
      \texttt{HANDSHAKE\_FINISH} message:
      \[
            \{version,\ session\_id,\ ct_A\}.
      \]
      Bob decapsulates:
      \[
            ss_A = \text{ML-KEM.Decaps}(ct_A, EK^{kem}_{priv,B}).
      \]
      Both parties now possess the same pair of shared secrets
      $(ss_A, ss_B)$ and derive a common master secret:
      \[
            master\_secret = \text{HKDF}(ss_A \parallel ss_B \parallel H(transcript)).
      \]
      Because the transcript hash $H(transcript)$ is included, the derived key
      is cryptographically bound to the exact handshake that was authenticated.

      \subsubsection{Step 5: Session key schedule and authentication state}
      From the master secret, Alice and Bob derive a root key and two symmetric
      chain keys:
      \begin{align*}
            root\_key      &= \text{HKDF}(master\_secret,\texttt{root}) \\
            CK_{A \rightarrow B} &= \text{HKDF}(root\_key,\texttt{A\rightarrow B}) \\
            CK_{B \rightarrow A} &= \text{HKDF}(root\_key,\texttt{B\rightarrow A})
      \end{align*}

      The session state stored by each client contains:
      \begin{itemize}
            \item the peer identity key fingerprint,
            \item the session identifier,
            \item the current root key,
            \item one sending chain key and one receiving chain key,
            \item message counters and ratchet counters.
      \end{itemize}

      If identity verification has already been done out of band, the peer
      fingerprint must match the pinned value. Otherwise, the first verified
      connection can be accepted in a trust-on-first-use model.

      \subsubsection{Step 6: Protected message exchange}
      For each outgoing message, the sender advances the sending chain:
      \begin{align*}
            msg\_key_i &= \text{HKDF}(CK,\texttt{msg} \parallel i) \\
            CK_{next}  &= \text{HKDF}(CK,\texttt{chain} \parallel i)
      \end{align*}
      The plaintext is then encrypted with AES-256-GCM:
      \[
            C_i = \text{AES-256-GCM}(msg\_key_i,\ nonce_i,\ plaintext_i,\ AAD_i).
      \]
      The additional authenticated data $AAD_i$ should contain metadata that
      must be integrity-protected but not encrypted, for example:
      \[
            AAD_i = session\_id \parallel sender\_id \parallel counter \parallel ratchet\_counter \parallel version.
      \]
      After a message key is used, it must be securely erased. This provides a
      simplified form of forward secrecy at the message layer.

      \subsubsection{Step 7: Simplified post-quantum ratchet}
      A full Signal-style Double Ratchet can be complex to implement. For this
      project, a simplified ratchet is sufficient:
      \begin{enumerate}
            \item \textbf{Symmetric ratchet:} derive a new message key from the
                  current chain key for every message.
            \item \textbf{Periodic PQ rekeying:} every $N$ messages, or after a
                  fixed timeout, the sender generates a fresh ephemeral ML-KEM
                  public key and signs it with ML-DSA.
            \item \textbf{Root key refresh:} the receiver encapsulates to that
                  fresh public key, both sides derive a new shared secret, and
                  mix it into the root key with HKDF.
      \end{enumerate}
      This periodic rekeying improves post-handshake forward secrecy without
      requiring the full complexity of the Double Ratchet protocol.

      \subsubsection{Security properties}
      The protocol aims to provide the following guarantees:
      \begin{itemize}
            \item \textbf{Confidentiality:} messages are encrypted with
                  AES-256-GCM under per-message keys.
            \item \textbf{Authentication:} only a party owning the correct
                  ML-DSA identity key can produce valid handshake signatures.
            \item \textbf{Forward secrecy:} ephemeral ML-KEM private keys are
                  deleted after use, and message keys are not reused.
            \item \textbf{Man-in-the-middle resistance:} an attacker cannot
                  replace authenticated keys without forging ML-DSA signatures.
            \item \textbf{Quantum resistance:} the public-key layer does not
                  rely on ECC or RSA, avoiding the harvest-now-decrypt-later
                  weakness of classical public-key protocols.
      \end{itemize}

      \subsubsection{MITM simulation plan}
      To evaluate the practical security implications, the prototype can be
      tested under three scenarios:
      \begin{enumerate}
            \item \textbf{Unauthenticated exchange:} the attacker intercepts and
                  replaces public keys. This should succeed in both the classical
                  and post-quantum versions because authentication is absent.
            \item \textbf{Authenticated classical protocol:} the attacker can
                  capture the entire transcript, including public keys,
                  ciphertexts, signatures, and metadata, but cannot tamper with
                  the session without being detected. However, recorded traffic
                  remains vulnerable to future quantum attacks against X25519 or
                  Ed25519.
            \item \textbf{Authenticated post-quantum protocol:} the attacker can
                  still capture the full transcript, but the confidentiality of
                  recorded sessions should remain protected against known quantum
                  attacks because the protocol uses ML-KEM and ML-DSA.
      \end{enumerate}

      \subsubsection{Comparative benchmark plan}
      The benchmark should compare a classical profile
      (X25519 + Ed25519 + AES-256-GCM) with a post-quantum profile
      (ML-KEM + ML-DSA + AES-256-GCM). The following metrics should be measured:
      \begin{itemize}
            \item public key size,
            \item private key size,
            \item KEM ciphertext size,
            \item signature size,
            \item total handshake bytes exchanged,
            \item handshake latency,
            \item encapsulation and decapsulation time,
            \item signing and verification time.
      \end{itemize}

      In practice, the expected result is that the post-quantum protocol offers
      stronger long-term security against quantum adversaries, but at the cost of
      significantly larger keys, larger signatures, and higher handshake
      overhead.

      \section{Some examples to get started}

      \subsection{How to include Figures}

      First you have to upload the image file from your computer using the
      upload link in the file-tree menu. Then use the includegraphics command to
      include it in your document. Use the figure environment and the caption
      command to add a number and a caption to your figure. See the code for Figure
      \ref{fig:frog} in this section for an example.

      Note that your figure will automatically be placed in the most appropriate
      place for it, given the surrounding text and taking into account other
      figures or tables that may be close by. You can find out more about adding
      images to your documents in this help article on \href{https://www.overleaf.com/learn/how-to/Including_images_on_Overleaf}{including
      images on Overleaf}.

      \begin{figure}
            \centering
            \includegraphics[width=0.25\linewidth]{frog.jpg}
            \caption{\label{fig:frog}This frog was uploaded via the file-tree menu.}
      \end{figure}

      \subsection{How to add Tables}

      Use the table and tabular environments for basic tables --- see Table~\ref{tab:widgets},
      for example. For more information, please see this help article on
      \href{https://www.overleaf.com/learn/latex/tables}{tables}.

      \begin{table}
            \centering
            \begin{tabular}{l|r}
                  Item    & Quantity \\
                  \hline
                  Widgets & 42       \\
                  Gadgets & 13
            \end{tabular}
            \caption{\label{tab:widgets}An example table.}
      \end{table}

      \subsection{How to add Comments and Track Changes}

      Comments can be added to your project by highlighting some text and
      clicking ``Add comment'' in the top right of the editor pane. To view existing
      comments, click on the Review menu in the toolbar above. To reply to a comment,
      click on the Reply button in the lower right corner of the comment. You can
      close the Review pane by clicking its name on the toolbar when you're done
      reviewing for the time being.

      Track changes are available on all our \href{https://www.overleaf.com/user/subscription/plans}{premium
      plans}, and can be toggled on or off using the option at the top of the
      Review pane. Track changes allow you to keep track of every change made to
      the document, along with the person making the change.

      \subsection{How to add Lists}

      You can make lists with automatic numbering \dots

      \begin{enumerate}
            \item Like this,

            \item and like this.
      \end{enumerate}
      \dots or bullet points \dots
      \begin{itemize}
            \item Like this,

            \item and like this.
      \end{itemize}

      \subsection{How to write Mathematics}

      \LaTeX{} is great at typesetting mathematics. Let
      $X_{1}, X_{2}, \ldots, X_{n}$ be a sequence of independent and identically
      distributed random variables with $\text{E}[X_{i}] = \mu$ and
      $\text{Var}[X_{i}] = \sigma^{2}< \infty$, and let
      \[
            S_{n}= \frac{X_{1}+ X_{2}+ \cdots + X_{n}}{n}= \frac{1}{n}\sum_{i}^{n}
            X_{i}
      \]
      denote their mean. Then as $n$ approaches infinity, the random variables $\sqrt{n}
      (S_{n}- \mu)$ converge in distribution to a normal $\mathcal{N}(0, \sigma^{2}
      )$.

      \subsection{How to change the margins and paper size}

      Usually the template you're using will have the page margins and paper
      size set correctly for that use-case. For example, if you're using a
      journal article template provided by the journal publisher, that template will
      be formatted according to their requirements. In these cases, it's best not
      to alter the margins directly.

      If however you're using a more general template, such as this one, and would
      like to alter the margins, a common way to do so is via the geometry
      package. You can find the geometry package loaded in the preamble at the top
      of this example file, and if you'd like to learn more about how to adjust
      the settings, please visit this help article on \href{https://www.overleaf.com/learn/latex/page_size_and_margins}{page
      size and margins}.

      \subsection{How to change the document language and spell check settings}

      Overleaf supports many different languages, including multiple different languages
      within one document.

      To configure the document language, simply edit the option provided to the
      babel package in the preamble at the top of this example project. To learn
      more about the different options, please visit this help article on
      \href{https://www.overleaf.com/learn/latex/International_language_support}{international
      language support}.

      To change the spell check language, simply open the Overleaf menu at the top
      left of the editor window, scroll down to the spell check setting, and adjust
      accordingly.

      \subsection{How to add Citations and a References List}

      You can simply upload a \verb|.bib| file containing your BibTeX entries, created
      with a tool such as JabRef. You can then cite entries from it, like this:
      \cite{greenwade93}. Just remember to specify a bibliography style, as well
      as the filename of the \verb|.bib|. You can find a \href{https://www.overleaf.com/help/97-how-to-include-a-bibliography-using-bibtex}{video
      tutorial here} to learn more about BibTeX.

      If you have an \href{https://www.overleaf.com/user/subscription/plans}{upgraded
      account}, you can also import your Mendeley or Zotero library directly as
      a \verb|.bib| file, via the upload menu in the file-tree.

      \subsection{Good luck!}

      We hope you find Overleaf useful, and do take a look at our \href{https://www.overleaf.com/learn}{help
      library} for more tutorials and user guides! Please also let us know if
      you have any feedback using the \textbf{Contact us} link at the bottom of the
      Overleaf menu --- or use the contact form at \url{https://www.overleaf.com/contact}.

      \bibliographystyle{alpha}
      \bibliography{sample}
\end{document}
